\documentclass[11pt]{article}
\usepackage[margin=1in]{geometry}
\usepackage{amsmath,amssymb}
\newcommand{\finalanswer}[1]{\par\medskip\noindent\fbox{\parbox{0.94\linewidth}{\textbf{Answer:} #1}}}
\title{STAT 412 Homework 3 --- Question 4}
\author{}
\date{}
\begin{document}
\maketitle

\section*{Problem}
The density function of the continuous random variable $X$, the total number
of hours, in units of 100 hours, that a family runs a vacuum cleaner over a
period of one year, is given. Find the average number of hours per year that
families run their vacuum cleaners.
\[
f(x)=
\begin{cases}
\dfrac95x,&0<x<\dfrac{\sqrt5}{3},\\[3pt]
\dfrac{6\sqrt5}{5}-\dfrac95x,
&\dfrac{\sqrt5}{3}\le x<\dfrac{2\sqrt5}{3},\\[3pt]
0,&\text{elsewhere}.
\end{cases}
\]

The average number of hours per year that families run their vacuum cleaners is
\fbox{\strut\hspace{1.5em}} hours.

\textit{(Round to the nearest whole number as needed.)}

\section*{Work}
This triangular density is symmetric about its peak at
$x=\sqrt5/3$, so its mean is the center of symmetry:
\[
E(X)=\frac{\sqrt5}{3}.
\]
Since $X$ is measured in hundreds of hours,
\[
100E(X)=\frac{100\sqrt5}{3}\approx74.536\text{ hours}.
\]

\finalanswer{The average annual use is $\boxed{75\text{ hours}}$.}
\end{document}
